\scanpage{12}
si $s'$ est son point fermé, on trouve aussi un homomorphisme de \unsure{restriction}
\diagram{\begin{tikzcd}[ampersand replacement=\&](2)\qquad K_{\mathrm{BT}}(S',G)\arrow[r,dash]\&K_{\mathrm{BT}}(s',G),\end{tikzcd}}
en notant que $G$ opère sur $s'$ (de façon pas nécessairement fidèle). Or, $k'=k(s')$ étant supposé \deleted{complet}\inserted{parfait}, on a un foncteur \uline{covariant} naturel
\diagram{\begin{tikzcd}[ampersand replacement=\&]\mathrm{BT}(s')\deleted{\unsure{G}}\arrow[r,dash]\&\operatorname{Modlib}(W(k')),\end{tikzcd}}
\margin{et $=$ vers $\operatorname{Cris}(k')$}
et comme $G$ opère sur $k'$, donc sur $W(k')$, le foncteur précédent induit trivialement un foncteur
\diagram{\begin{tikzcd}[ampersand replacement=\&]\mathrm{BT}(s',G)\arrow[r]\&\operatorname{Modlib}(W(k'),G),\end{tikzcd}}
d'où par suite un homomorphisme
\diagram{\begin{tikzcd}[ampersand replacement=\&](3)\qquad K_{\mathrm{BT}}(s',G)\arrow[r,dash]\&K(\deleted{\unsure{\mathbf Z}}W(k'),G).\end{tikzcd}}
Composant (1),(2) et (3), on trouve un homomorphisme canonique
\diagram{\begin{tikzcd}[ampersand replacement=\&](4)\qquad K_{\mathrm{BT}}(K;K')\arrow[r]\&K(W(k'),G).\end{tikzcd}}
Notons d'ailleurs que si $k$ est algébriquement clos, alors $G$ opère trivialement sur $k'$ donc sur $W(k')$, et on a alors
\[
K(W(k'),G)=R_{W(k')}(G);
\]
dans le cas général, il y a lieu d'introduire le groupe d'intertie $I\subset G$, et on trouve un homomorphisme
\diagram{\begin{tikzcd}[ampersand replacement=\&]
K(W(k'),G)\arrow[r]\&K(W(\unsure{k'}),I)^{\inserted{\unsure{G|G/I}}}\simeq R_{W(k')}(I)^{\inserted{G/I}},
\end{tikzcd}}
\margin{\uline{NB} $G\deleted{\unsure{/I}}$ opère sur $R$ via ses opérations sur $I$ \uline{et} sur $W(k')$ par transport de structure}
\margin{On peut \unsure{dire en|dire ça}\ \unsure{réunissant} les différents points ....}
d'où grâce à (4)
\diagram{\begin{tikzcd}[ampersand replacement=\&]
(5)\qquad K_{\mathrm{BT}}(K,K')\arrow[r]\&R_{W(k')}(I)^{\inserted{\unsure{G/I|G}}}\arrow[r]\&\operatorname{Cent}_{\inserted{G}}(I,W(k')),
\end{tikzcd}}
la dernière flèche étant la flèche trace. L'homomorphisme composé est appelé encore l'homomorphisme trace, l'image de $\operatorname{cl}(M)$ dans le dernier membre est appelé le \uline{caractère} du groupe de B.T. $M$. On voit facilement que dans un sens évident, il ne dépend pas du choix du corps de bonne réduction virtuelle $K'$: la fonction caractère est alors envisagée comme fonction sur le \inserted{sous-}groupe d'inertie (infini) de $\operatorname{Gal}(\overline K/K)$.

Supposons que $M_{K'}$ soit \uline{constant}, de sorte que \deleted{\unsure{$M$ est un}\ \illegible{2}} $T_p(M)$ est un faisceau p-adique localement constant (au sens technique
\mspage{11}
